\section{記述統計量}
\label{lesson03}

\subsection{授業内容}
\subsubsection{科目の中でのこのコマの位置づけ}

\begin{tabular}[t]{ccccccc}
    記述統計[2] & 確率[0] & モデル[0] & 推測・推定[0] & 意思決定[0] & 実践[0]
\end{tabular}


\subsubsection{概要}
統計の基本は集めたデータを集計し，要約統計量で表現することである。手元にあるデータの詳細を記述するため，これは記述統計(descriptive statistics)と呼ばれる。いかなる進んだ分析を行うに当たっても，必ず可視化および記述統計による特徴の把握を欠かすことはできない。また要約統計量はその性質上，必ず有意味な情報を捨て去って代表的な性質を表すものであり，単一の指標だけで考えることは危険であることを理解し，活用できるようにする。
\subsubsection{コマ主題細目(3-5項目箇条書き)}
\begin{description}
\item[中心化傾向の指標] 散らばりを持つデータの中心についての要約統計量として，平均値(mean)，中央値(median) ,最頻値(mode)の三つがあげられる。それぞれの特徴および算出方法，その際の注意点についても理解する。また，計算や表現の利便性のため，$\sum$という総和記号を用いる。
\begin{flushright}
    $\to$ \citet{山内201001}のP.27--35,\citet{Yamada2004}のP.30--33，\citet{川端201408}のP.23--26.
\end{flushright}

\item[散らばりの指標] データにとって重要なのは中心化傾向の指標だけではなく，むしろその散らばりについての指標である分散(variance),標準偏差(standard deviation)こそ重要である。加えて最大値(max)，最小値(min) ,四分位偏差(quantile)，パーセンタイル(percentile)，範囲(range)などが散らばりを表す指標として用いられる。ここではそれぞれの特徴および算出法について理解する。
\begin{flushright}
    $\to$ \citet{山内201001}のP.37--48,\citet{Yamada2004}のP.34--37，\citet{川端201408}のP.26-33.
\end{flushright}


\item[データの相対的位置]平均値と標準偏差を導入したことで，データ点の相対的位置を表現することが可能になった。心理学の研究対象の多くは，絶対ゼロ点をもたない間隔尺度水準であり，個々のデータ点の特徴を記述する際には相対的な位置である標準得点で表されることが多い。ここでは標準化の方法と標準得点の特徴，標準得点の応用である偏差値の特徴および算出方法について理解する。
\begin{flushright}
    $\to$ \citet{山内201001}のP.54--56,\citet{Yamada2004}のP.38--41，\citet{川端201408}のP.36--40.
\end{flushright}

\item[変数間関係の指標]ここまでは一つの変数についての特徴を記述するものであったが，次に複数の変数による変数同士の関係を考える。一つは共分散であり，それを標準化したスコアにしたものが相関係数であることを理解する。また相関係数はあくまでも変数間の直線的関係についての記述であり，相関係数だけを見て判断すると対象の理解を誤る場合があることに注意する。
\begin{flushright}
    $\to$ \citet{山内201001}のP.62--75,\citet{Yamada2004}のP.44--59，\citet{川端201408}のP.48--58.
\end{flushright}


\end{description}
\subsubsection{キーワード}
\begin{itemize}
\item 平均値，中央値，最頻値
\item 分散，標準偏差，最大値，最小値，四分位偏差，パーセンタイル，範囲
\item 標準得点，偏差値
\item 共分散，相関係数
\end{itemize}
\subsection{授業情報}
\paragraph{コマの展開方法}
講義
\paragraph{標準シラバスにおける位置づけ}
科目番号5 心理学統計法；1.心理学で用いられる統計手法；A 記述統計:代表値と散布度/B 記述統計：相関
\subsubsection{予習・復習課題}
\paragraph{予習}
統計分野は2012年(平成24年)より，高校数学の必修である数学Iで新たに扱われることとなっているため，受講生諸君はすでに算術平均や相関係数について一度は習っていることと思われる。内容について自信がない場合は，高校数学のテキストを読み返し，キーワードにあげられている要約統計量について予習をしておくこと。また心理統計は統計学のユーザでもあるので，PCを使ってこれらの指標を計算できるように準備，練習しておくことが望ましい。
\paragraph{復習}
今回の内容については，心理統計の各種テキストの中に必ず含まれている内容であるから，類書の当該箇所について自分なりの復習をしてもらいたい。\citet{Yamada2004}はユーザー目線から丁寧な記述がされている良書であり，特に以降の推測統計については慎重かつ緻密な言葉遣いで記述されていることから，日本の多くの心理統計業界でのテキストとして愛用されている。\citet{川端201408}は\citet{Yamada2004}以上に丁寧かつ細部にまで言及され，言葉遣いも柔らかいので大変読みやすく感じる。特にこの本はシリーズ化されていて，研究領域ごとに特化した統計分析がテーマになっている\footnote{なお，このシリーズの各巻末にある著者紹介は大変ユニークであるので是非順番に読んでみよう。}。\citet{山内201001}は数学的記号が多用されており，敷居が高く感じる読者もいるかもしれないが，一冊で網羅的に記述されており相関・回帰を先に導入するなどモデルベースの観点から記述されている点が特徴的えだる。

数学的表現は難解に思えるかもしれないが，必要十分な情報が全て盛り込まれているし，何より積み重ねの学問であるので，不明な点が出た場合は定義に戻って確認しておくことを習慣付けよう。また記号だけでの理解ではなく，実際に数点のデータをとって，各統計量を手計算，あるいはPCを使って算出する練習をしておくと良い。記号$\sum$に関する例題としては，\citet{山内201001}のP.10--13が適当である。
